\question{How can the budget be reduced without losing accuracy?}
\label{sec:reduce}
\statusMeasured

\begin{shortanswer}
The cheapest lever is the round count, and it is free. The logistic protocols are fully converged by
round~10, but the configuration ran to \DPRounds{} and beyond; every round composes additional privacy
loss, so the cheap cut here is 50$\to$10 rounds: the repository's own R\'enyi
accountant at $\sigma = 2.0$, $\delta = 10^{-5}$ prices 50 rounds at $\varepsilon = 22.0$ against
8.1 at ten (a 2.7$\times$ reduction -- composition is sublinear, not linear) at \emph{no measured
accuracy cost}.
That change is now the default ($R = 10$ rounds). The second lever is now landed too: the record-level
DP-SGD standard (\S\ref{sec:noise}) orchestrates each practice's batch size and noise multiplier
\emph{individually}, so the smaller practices' budgets no longer overrun under a flat configuration,
under the toy schedule at ten local epochs a
uniform $(b, \sigma)$ delivers a composed $\varepsilon$ of 1.10 for a 50,000-patient practice and
16.04 for a 500-patient one; on this report's ten-practice census under the deployment schedule
(two epochs, ten rounds) the same flat instruction already spans 7.25--15.26 -- while the orchestrator contracts every
practice to one shared target $\varepsilon^{*}$. Beyond these, the levers are more practices the
pilot's twenty-five rather than today's \NumPractices{} and subsampling amplification, which the
standard already earns in part through its Poisson lots.
\end{shortanswer}

\subsection{Convergence is reached long before the round budget is spent}

\begin{figure}[htbp]
  \centering
  \includegraphics[width=\textwidth]{convergence}
  \caption{AUROC by round, \NumSeeds{} seeds, shaded $\pm$ 1 s.d. The logistic protocols are flat
  from round~10 onward. The tree-based protocol \emph{degrades} with more rounds. FedAvg is not
  missing it lies underneath FedProx, which is the null result reported in
  \S\ref{sec:modelchoice}: the two agree to three decimals.}
  \label{fig:convergence}
\end{figure}

\begin{table}[htbp]
  \centering
  \caption{Accuracy by round count. Nothing after round~10 is bought with the extra privacy loss.}
  \setlength{\tabcolsep}{4.5pt}\small
\begin{tabular}{rrrrrrrrrrr}
    \toprule
    & \multicolumn{5}{c}{Global AUROC} & \multicolumn{5}{c}{Worst practice} \\
    \cmidrule(lr){2-6} \cmidrule(lr){7-11}
    Rounds & Prox & XGB & Mosaic & Local & Avg & Prox & XGB & Mosaic & Local & Avg \\
    \midrule
    5 & 0.805 & 0.780 & 0.800 & 0.791 & 0.805 & 0.672 & 0.599 & 0.655 & 0.511 & 0.671 \\
    10 & 0.807 & 0.757 & 0.800 & 0.790 & 0.807 & 0.675 & 0.600 & 0.663 & 0.498 & 0.676 \\
    20 & 0.808 & 0.729 & 0.800 & 0.789 & 0.808 & 0.675 & 0.555 & 0.661 & 0.495 & 0.675 \\
    35 & 0.808 & 0.703 & 0.800 & 0.786 & 0.808 & 0.677 & 0.521 & 0.653 & 0.462 & 0.677 \\
    50 & 0.808 & 0.686 & 0.799 & 0.785 & 0.808 & 0.677 & 0.528 & 0.646 & 0.462 & 0.677 \\
    \bottomrule
\end{tabular}

  \label{tab:rounds}
\end{table}

Reading the table across rows rather than down columns is the point: for every logistic comparator,
the round-10 figure and the round-\MaxRounds{} figure agree to three decimals. Those forty extra rounds
composed forty extra rounds of privacy loss and returned nothing.

\textbf{The deployment default is now 10 rounds.} At $\sigma = 2.0$ this takes the composed budget
from \EpsAtSigmaTwo{} to approximately 8.1 the largest single improvement available, at zero
accuracy cost and zero engineering risk.

One protocol behaves differently and it is worth noting here because it interacts with the same
lever: the tree-based protocol peaks early and then degrades monotonically, because each round
appends every practice's newly grown trees to one ensemble and, under non-IID data, those trees
encode contradictory local rules. More rounds means more contradiction. For that protocol the round
count is a hyperparameter to tune \emph{down}, not up but \S\ref{sec:modelchoice} recommends
against deploying it at all.

\subsection{The other levers, in order of effect}

\begin{enumerate}
  \item \textbf{Per-practice orchestration of $(b, \sigma)$ landed.} Under any flat DP-SGD
        configuration the composed budget is an accident of practice size: identical settings of
        batch~64 and $\sigma = 1.5$ deliver $\varepsilon = 1.10$ to a 50,000-patient practice and
        $\varepsilon = 16.04$ to a 500-patient one (toy endpoints, ten local epochs; the
        deployment-schedule spread across this report's ten-practice census is 7.25--15.26). The standard's rule-based server planner inverts the shared R\'enyi-DP accountant
        per practice (\S\ref{sec:noise}), so one target $\varepsilon^{*}$ is certified for every
        practice at $\varepsilon^{*} = 2$ the planned $\sigma$ spans 8.32--11.68, and the measured
        utility is flat from $\varepsilon^{*} = 0.5$ to $8$. What used to be ``reduce the budget
        where it hurts'' becomes ``choose the budget once, for everyone''.
  \item \textbf{More practices.} The budget is driven by how much population an individual can hide
        in. Twenty-five practices instead of \NumPractices{} is the pilot's own plan and is the
        structural fix. Under orchestration its effect shows up directly as a smaller required
        $\sigma$ per practice rather than as an easier target.
  \item \textbf{Subsampling amplification.} The standard already credits amplification
        \emph{within} each practice: its Poisson lots sample records at rate $q = b/N_k$, and the
        accountant prices $q$ explicitly. Remaining on the table is amplification \emph{across}
        practices if a practice does not participate in every round, its exposure is further
        amplified away in the accountant's favour. The current configuration has every practice
        train every round, which earns no such credit; the cost of taking it is slower convergence,
        which the round-count finding above suggests there is room to absorb.
  \item \textbf{Accept a larger $\varepsilon$ and let secure aggregation carry the load.} Secure
        aggregation costs bandwidth and coordination, not accuracy (\S\ref{sec:secagg}). At this
        cohort size it protects the individual contribution far more cheaply than noise does, and
        the two are complementary rather than alternatives and they compose in one run:
        \texttt{dpsgd-epsilon} with \texttt{secure-aggregation=true} has been executed live with
        zero failures.
\end{enumerate}

\subsection{What the binding constraint really is}

Throughout, the constraint that binds first is not the mean but the \emph{worst practice}. Global
AUROC degrades gently with noise; the worst practice degrades sharply and with much larger variance.
Since the project's premise is that small practices should get a model they could not build alone,
the worst-practice figure is the one that decides whether a privacy setting is deployable and it
is also the T2.5 fairness measure. Any budget proposal should be evaluated against that column, not
against the headline.
